\chapter{小数}\label{ch:g6:decimals}

在 $2$ 与 $3$ 之间有许多数：$2.5$、$2.71$、
$2.999$\,\dots\ 小数把数位系统推广到
个位的右侧，它们与整数遵循相同的规则——
本章将教你避开其中的几个陷阱。

\section{小数的写法}

\begin{definition}[小数的数位]\label{def:g6:decimals:places}
一个\emph{小数}\index{小数}由整数部分和
小数部分组成，二者以小数点隔开。小数点右边的每一个数位
的值都是前一位的十分之一：十分位、
百分位、千分位。例如
\[
13.407 = 13 + \frac{4}{10} + \frac{0}{100} + \frac{7}{1000}
= 1 \times 10 + 3 + 4 \times \frac{1}{10} + 7 \times \frac{1}{1000}.
\]
\end{definition}

\begin{omfigure}
\begin{tikzpicture}[scale=0.9]
  \draw (0,0) grid[xstep=2.3, ystep=0.85] (11.5,1.7);
  \node[font=\small] at (1.15,1.275) {十位};
  \node[font=\small] at (3.45,1.275) {个位};
  \node[font=\small] at (5.75,1.275) {十分位};
  \node[font=\small] at (8.05,1.275) {百分位};
  \node[font=\small] at (10.35,1.275) {千分位};
  \node[font=\large, omDef] at (1.15,0.425) {$1$};
  \node[font=\large, omDef] at (3.45,0.425) {$3$};
  \node[font=\large, omThm] at (5.75,0.425) {$4$};
  \node[font=\large, omThm] at (8.05,0.425) {$0$};
  \node[font=\large, omThm] at (10.35,0.425) {$7$};
  \fill (4.6,0.06) circle (2.5pt);
\end{tikzpicture}

{\small 数 $13.407$ 的数位表：小数点位于
个位与十分位之间。}
\end{omfigure}

\begin{remark}[有意义的零与无意义的零]
在小数部分的\emph{末尾}添零不改变数值：
$2.5 = 2.50 = 2.500$。但数字\emph{之间}的零是必不可少的：
$13.407 \neq 13.47$。而 $0.5$ 与 $0.05$ 大不相同！
\end{remark}

\begin{method}[比较小数的大小]\label{met:g6:decimals:compare}
\begin{enumerate}
  \item 先比较整数部分：$7.2 > 5.99$，因为 $7 > 5$。
  \item 若整数部分相等，则逐位比较小数部分，
        从左往右比较 --- 在末尾补零有助于比较：
        比较 $3.4$ 与 $3.15$ 时，写成 $3.40$ 与 $3.15$；
        由于 $40 > 15$（以百分位计），所以 $3.4 > 3.15$。
\end{enumerate}
注意：\emph{位数多未必数值大}。$3.15$ 的位数比
$3.4$ 多，但数值更小。
\end{method}

\begin{omfigure}
\begin{tikzpicture}[scale=1.1]
  \draw[-Stealth] (-0.3,0) -- (10.6,0);
  \foreach \i in {0,1,...,10} \draw (\i,-0.08) -- (\i,0.08);
  \node[below, font=\small] at (0,-0.1) {$3$};
  \node[below, font=\small] at (5,-0.1) {$3.5$};
  \node[below, font=\small] at (10,-0.1) {$4$};
  \fill[omThm] (1.5,0) circle (2pt) node[above, font=\small] {$3.15$};
  \fill[omDef] (4,0) circle (2pt) node[above, font=\small] {$3.4$};
  \fill[omMeth] (9,0) circle (2pt) node[above, font=\small] {$3.9$};
\end{tikzpicture}

{\small 放大看数轴上 $3$ 与 $4$ 之间的部分：每个刻度代表
一个十分位。点 $3.15$ 位于 $3.1$ 与 $3.2$ 的正中间。}
\end{omfigure}

\section{小数的加法、减法与乘法}

\begin{method}[小数的竖式计算]\label{met:g6:decimals:columns}
进行小数的加法或减法时，先将\emph{小数点}对齐
（即个位对个位、十分位对十分位），必要时在末尾补零；
再像整数那样计算，并使结果的小数点
与上面各数的小数点对齐。
\end{method}

\begin{example}\label{ex:g6:decimals:addsub}
计算 $13.7 + 2.85$，对齐小数点（$13.7 = 13.70$）：
\[
\begin{array}{r}
1\,3.7\,0 \\
+\ \ \,2.8\,5 \\
\hline
1\,6.5\,5
\end{array}
\]
计算 $6 - 2.35$（写成 $6 = 6.00$）：
$6.00 - 2.35 = 3.65$。验算：$3.65 + 2.35 = 6$。
\end{example}

\begin{example}[小数的乘法]\label{ex:g6:decimals:mult}
计算 $2.3 \times 1.4$：先把两个因数当作整数相乘，
即 $23 \times 14 = 322$；再数出两个因数小数部分的位数
（$1 + 1 = 2$），点上小数点，使结果恰有这么多位小数：
$2.3 \times 1.4 = 3.22$。估算检验：$2.3 \times 1.4$
应略大于 $2.3 \times 1 = 2.3$ --- 而 $3.22$ 确实如此。
\end{example}

\section{乘以与除以 10、100、1000}

\begin{proposition}[移动小数点]\label{prop:g6:decimals:shift}
一个小数乘 $10$、$100$、$1000$ 时，其小数点
分别向\emph{右}移动 $1$、$2$、$3$ 位；除以它们时，
则向\emph{左}移动（必要时补零）：
\[
3.75 \times 100 = 375, \qquad
42.1 \div 1000 = 0.0421 .
\]
\end{proposition}

\begin{proof}
乘以 $10$ 使每个数字的值变为原来的十倍：十分位
变为个位，个位变为十位，依此类推——每个数字
在数位表中向左移动一列，这相当于把小数点向右移动一位。
除法则作相反的变化。
\end{proof}

\begin{example}[单位换算]\label{ex:g6:decimals:units}
单位换算正是这一游戏：由于 $1$ m $= 100$ cm，
\[
3.75 \text{ m} = 3.75 \times 100 \text{ cm} = 375 \text{ cm},
\qquad
42 \text{ mm} = 42 \div 10 \text{ cm} = 4.2 \text{ cm}.
\]
\end{example}

\section{四舍五入}

\begin{definition}[四舍五入]\label{def:g6:decimals:rounding}
一个数精确到个位（或十分位、
百分位、\dots）的\emph{四舍五入}\index{四舍五入}是与其最接近的、具有该精度的数。看
下一位数字：若为 $0,1,2,3,4$，则舍去（保留）；若为
$5,6,7,8,9$，则进位。
\end{definition}

\begin{example}
$7.38$ 精确到个位为 $7$（下一位为 $3$，舍去）；精确到
十分位则为 $7.4$（下一位为 $8$，进位）。价格
$4.996$ 精确到百分位为 $5.00$——四舍五入可以改变
每一位数字！
\end{example}

\section{练习}

\begin{exercise}[$\star$]\label{exo:g6:decimals:1}
将下列各数写成小数：$5 + \dfrac{3}{10} + \dfrac{7}{100}$；
\quad $\dfrac{9}{10}$; \quad $12 + \dfrac{4}{1000}$;
\quad “八个一与五个百分之一”。
\end{exercise}

\begin{exercise}[$\star$]\label{exo:g6:decimals:2}
在 $86.354$ 中：十分位上的数字是多少？百分位上的数字呢？
数字 $8$ 计数的单位是什么？请按照
\cref{def:g6:decimals:places}中的形式写出这个数。
\end{exercise}

\begin{exercise}[$\star$]\label{exo:g6:decimals:3}
抄写并在问号处填入 $<$、$>$ 或 $=$：
\[
5.3 \;?\; 5.29, \qquad
0.7 \;?\; 0.70, \qquad
2.09 \;?\; 2.9, \qquad
14.5 \;?\; 14.49 .
\]
\end{exercise}

\begin{exercise}[$\star$]\label{exo:g6:decimals:4}
按从小到大的顺序排列：
$4.2$; \ $4.05$; \ $4.51$; \ $4.15$; \ $4.5$.
\end{exercise}

\begin{exercise}[$\star$]\label{exo:g6:decimals:5}
在以十分之一为刻度的数轴上，点 $A$、$B$、$C$ 分别对应哪些小数，
已知 $A$ 位于 $3$ 格处（起点为 $6$），$B$ 位于 $7$ 格处（起点为 $6$），
$C$ 位于 $2$ 格处（起点为 $7$）？
\end{exercise}

\begin{exercise}[$\star$]\label{exo:g6:decimals:6}
用竖式计算：$45.8 + 7.65$；\quad $23.4 - 8.72$；\quad
$5.6 \times 2.4$（注意数一数小数的位数！）。
\end{exercise}

\begin{exercise}[$\star$]\label{exo:g6:decimals:7}
不列竖式，直接口算：
\[
6.42 \times 10, \qquad
0.35 \times 1000, \qquad
78.1 \div 100, \qquad
5 \div 1000 .
\]
\end{exercise}

\begin{exercise}[$\star$]\label{exo:g6:decimals:8}
换算：$2.4$ m 化为 cm；\quad $370$ g 化为 kg；\quad $0.85$ km 化为
m；\quad $56$ mm 化为 m。
\end{exercise}

\begin{exercise}[$\star$]\label{exo:g6:decimals:9}
对 $23.867$ 四舍五入：分别精确到个位、十分位和百分位；将
$9.97$ 四舍五入到十分位。
\end{exercise}

\begin{exercise}[$\star\star$]\label{exo:g6:decimals:10}
一根法棍 $1.15$ 元。莉娜买了三根法棍，她用一张
$5$ 元纸币付款。写出所需的两个算式，并求她应找回的钱数。
\end{exercise}

\begin{exercise}[$\star\star$]\label{exo:g6:decimals:11}
求一个严格介于 $7.4$ 与 $7.5$ 之间的小数；再求一个介于
$3.99$ 与 $4$ 之间的小数。两种情形下，这样的数各有多少个？
\end{exercise}

\begin{exercise}[$\star\star\star$]\label{exo:g6:decimals:12}
用数字 $2$、$5$、$8$ 各恰好用一次，并恰好使用一个
小数点，写出可能的最大数，再写出可能的最小数。
（形如 $.58$ 的数不允许，整数部分
至少要含有一位数字。）
\end{exercise}

\section{问题：3 之后的下一个数不存在}

\begin{problem}[{周末问题 --- 任意两个小数之间
总有另一个小数：放大数轴}]
\label{pb:g6:decimals:1}
在整数这一阶梯上，每个数都有一个紧挨着的后继数：
紧接在 $7$ 之后的是 $8$，两者之间
什么也没有。小数则是一个完全不同的世界。
$3$ 之后紧接着的那个数是什么？是 $3.1$ 吗？$3.01$ 吗？$3.001$ 吗？
本题在数轴上发展一种\emph{放大}技巧（承接
\cref{exo:g6:decimals:11}），并得出一个著名结论：
$3$ 之后的那个数\emph{不存在} ---
任意两个小数之间，无论它们多么接近，总还有
余地容纳更多的数。

\medskip
\noindent\textbf{第一部分 --- 放大观察。}
\begin{enumerate}
  \item 比较 $2.999$ 与 $3$ 的大小，
        并计算两者之差。
  \item 画一条从 $7.4$ 到 $7.5$ 的数轴，按百分位标出刻度，
        并在其上标出 $7.42$、$7.45$ 和 $7.48$。
        列出\emph{全部}两位小数，
        它们严格位于 $7.4$ 与 $7.5$ 之间。共有多少个？
  \item 再次放大：在 $7.44$ 与 $7.45$ 之间，列出所有
        三位小数的数。这一次
        共有多少个？
  \item 用同样的思路，说明如何计数
        恰好具有三位小数、且严格位于
        $7.4$ 与 $7.5$ 之间的那些数，并给出其个数。
  \item 归纳第 2--4 题背后隐藏的方法（即
        \emph{放大}）：给定两个看似相邻的数，
        如 $5.67$ 与 $5.68$，多写一位小数如何
        总能揭示出一个严格位于它们之间的数？
        将你的方法用于 $5.67$ 与 $5.68$，
        再用于 $0.1999$ 与 $0.2$。
\end{enumerate}

\medskip
\noindent\textbf{第二部分 --- 缺失的紧邻数。}
\begin{enumerate}[resume]
  \item 汤姆声称：“$3$ 之后的那个数是 $3.1$。”请证明
        他错了：指出一个严格位于 $3$ 与
        $3.1$ 之间的数。汤姆退而改说 $3.01$，再改说 $3.001$，请逐一驳倒
        他提出的每一个候选数。
  \item 说明为什么\emph{没有人}能在这个游戏中赢过你：
        无论汤姆提出哪个严格大于 $3$ 的数，
        第 5 题中的放大方法总能给出一个严格位于
        $3$ 与他所提的数之间的数。这说明了关于
        “$3$ 之后的那个数”的什么结论？
  \item 现在看另一侧：汤姆寻找
        \emph{紧接在} $3$ 之前的那个数，先后试了 $2.9$、$2.99$，
        再试 $2.999$。请逐一驳倒他的这三个候选数。然后
        计算 $3 - 2.999$，并（不用竖式计算）描述
        $3$ 与由数字 $2$、一个小数点和二十个
        数字 $9$ 所写成的数之间的差。
  \item 为什么上述方法对\emph{整数}行不通？
        请用一两句话说明：为什么没有整数
        严格位于 $7$ 与 $8$ 之间，尽管它们之间
        有很多小数。
  \item 某长度被报告为“$3.7$ cm，四舍五入到
        十分位”（见\cref{def:g6:decimals:rounding}）。求出
        能四舍五入得到 $3.7$ cm 的最小长度，并说明为什么
        长为 $3.75$ cm 的长度\emph{不会}四舍五入为 $3.7$ ---
        因此其真实长度至少为 $3.65$ cm，且严格
        小于 $3.75$ cm。
\end{enumerate}

\medskip
\noindent\textbf{第三部分 --- 数字与小数点的游戏。}
\begin{enumerate}[resume]
  \item 从小到大排列：$0.6$；\ $0.58$；\
        $0.123$；\ $0.0999$。然后用一句话向同学解释，
        为什么 $0.0999 < 0.6$，尽管它
        写出的位数更多（见\cref{met:g6:decimals:compare}）。
  \item 商店里的价格总是恰好有两位小数。
        是否存在严格位于 $4.99$ 与 $5$ 之间的\emph{价格}？
        请与第 9 题比较：价格与整数
        有什么共同之处？
  \item 用数字 $2$、$5$、$8$ 各恰好用一次，并恰好使用
        一个小数点（整数部分不能为空，见
        \cref{exo:g6:decimals:12}），求最接近
        $6$ 的数。通过计算你的最佳候选数与
        $6$ 的距离来说明理由。
  \item 数学中最著名的数 $\pi$ 满足
        $3.141 < \pi < 3.142$。请写出一个严格位于
        $3.141$ 与 $3.142$ 之间的小数；再写出一个严格位于
        $3.1415$ 与 $3.1416$ 之间的数。（数学家正是这样夹逼 $\pi$ ---
        每多写一位小数，就多做一次放大。
        你将在
        \cref{ch:g6:measure}中再次遇见 $\pi$。）
  \item 压轴题：说明为什么\emph{$0$ 与 $1$ 之间的小数
        个数比你能写出的任何数都多}。
        （一次放大能产生多少个数？放大过程能否
        永远停止？）
\end{enumerate}
\end{problem}
